% GENERATED FILE. Edit canonical lessons, then run npm run generate:books.
% canonical-source-hash: fnv1a64:d904d0c29c8f52a8
% canonical-lessons: PT-C380-ministro, PT-C380-partido, PT-C380-parlamento, PT-C380-eleger, PT-C380-eleitor

\chapter{Se eu fosse ministro}
\label{ch:pt-se-eu-fosse-ministro}
\hlchaptermodality{380}

\begin{chapteropening}
\textbf{By the end of this chapter:} \emph{I can say what I would do in politics: ministers, parties, parliament, elections and voters.}

Complete the last lesson of chapter 380: I can say what I would do in politics: ministers, parties, parliament, elections and voters.
\end{chapteropening}

\section[o ministro]{o ministro — the minister}

\label{lesson:PT-C380-ministro}



\begin{quote}
Before the new one: say the Portuguese for the average, then the Portuguese for the bachelor's degree.
\end{quote}

\subsection*{You'll want to know: o ministro}
\textbf{o ministro} — \textquotedblleft{}the minister\textquotedblright{}.

\begin{quote}
Se eu fosse ministro, construía mais escolas.
\end{quote}

\begin{quote}
\emph{If I were a minister, I would build more schools.}
\end{quote}

Feminine \emph{a ministra}. \emph{O primeiro-ministro} — the prime minister.

\begin{cousinweb}
\emph{ministro} comes from Latin \emph{minister} \textquotedblleft{}servant, attendant\textquotedblright{} (EtymDB), built on \emph{minus} \textquotedblleft{}less\textquotedblright{} (Lewis \& Short): first a servant of the state.
\end{cousinweb}

\begin{grammarlens}[title={What you've built}]
One more word for telling, explaining and describing at B1.
\end{grammarlens}

\subsection*{Your turn}
\begin{itemize}
\raggedright
  \item \emph{Say it:} \emph{o ministro}
  \item \emph{Say it:} \emph{o ministro}, once more
  \item \emph{Say it:} a sentence of your own with \emph{o ministro}
  \item \emph{Recall:} say the Portuguese for the average, then the Portuguese for the bachelor's degree, then say \emph{o ministro} again
\end{itemize}

\begin{tcolorbox}[breakable,colback=teal!4,colframe=teal!35!black,title={Before you move on}]
What does o ministro mean? (\textquotedblleft{}The minister\textquotedblright{}.) Say it once more.
\end{tcolorbox}
\section[o partido]{o partido — the political party}

\label{lesson:PT-C380-partido}



\begin{quote}
Before the new one: say the Portuguese for the bachelor's degree, then the Portuguese for the minister.
\end{quote}

\subsection*{You'll want to know: o partido}
\textbf{o partido} — \textquotedblleft{}the political party\textquotedblright{}.

\begin{quote}
Ele vota sempre no mesmo partido.
\end{quote}

\begin{quote}
\emph{He always votes for the same party.}
\end{quote}

Also the participle of \emph{partir}: \emph{o copo está partido} — the glass is broken.

\begin{cousinweb}
\emph{partido} is the participle of \emph{partir} \textquotedblleft{}to divide\textquotedblright{} used as a noun for a side or faction; EtymDB traces it through Old Portuguese \emph{partido} to \emph{partir}, which is inherited from Latin \emph{partīre} \textquotedblleft{}to share, to divide\textquotedblright{}, on \emph{pars} \textquotedblleft{}part\textquotedblright{} (EtymDB, Lewis \& Short).
\end{cousinweb}

\begin{grammarlens}[title={What you've built}]
One more word for telling, explaining and describing at B1.
\end{grammarlens}

\subsection*{Your turn}
\begin{itemize}
\raggedright
  \item \emph{Say it:} \emph{o partido}
  \item \emph{Say it:} \emph{o partido}, once more
  \item \emph{Say it:} a sentence of your own with \emph{o partido}
  \item \emph{Recall:} say the Portuguese for the bachelor's degree, then the Portuguese for the minister, then say \emph{o partido} again
\end{itemize}

\begin{tcolorbox}[breakable,colback=teal!4,colframe=teal!35!black,title={Before you move on}]
What does o partido mean? (\textquotedblleft{}The political party\textquotedblright{}.) Say it once more.
\end{tcolorbox}
\section[o parlamento]{o parlamento — the parliament}

\label{lesson:PT-C380-parlamento}



\begin{quote}
Before the new one: say the Portuguese for the minister, then the Portuguese for the political party.
\end{quote}

\subsection*{You'll want to know: o parlamento}
\textbf{o parlamento} — \textquotedblleft{}the parliament\textquotedblright{}.

\begin{quote}
O parlamento aprovou a lei ontem.
\end{quote}

\begin{quote}
\emph{Parliament passed the law yesterday.}
\end{quote}

Its members are \emph{os deputados}.

\begin{cousinweb}
\emph{parlamento} is usually traced to Old French \emph{parlement} \textquotedblleft{}discussion, meeting, assembly\textquotedblright{} (EtymDB), built on \emph{parler} \textquotedblleft{}to speak\textquotedblright{}, which goes back to Late Latin \emph{parabolāre}, on \emph{parabola} \textquotedblleft{}a comparison\textquotedblright{} (Lewis \& Short).
\end{cousinweb}

\begin{grammarlens}[title={What you've built}]
One more word for telling, explaining and describing at B1.
\end{grammarlens}

\subsection*{Your turn}
\begin{itemize}
\raggedright
  \item \emph{Say it:} \emph{o parlamento}
  \item \emph{Say it:} \emph{o parlamento}, once more
  \item \emph{Say it:} a sentence of your own with \emph{o parlamento}
  \item \emph{Recall:} say the Portuguese for the minister, then the Portuguese for the political party, then say \emph{o parlamento} again
\end{itemize}

\begin{tcolorbox}[breakable,colback=teal!4,colframe=teal!35!black,title={Before you move on}]
What does o parlamento mean? (\textquotedblleft{}The parliament\textquotedblright{}.) Say it once more.
\end{tcolorbox}
\section[eleger]{eleger — to elect}

\label{lesson:PT-C380-eleger}



\begin{quote}
Before the new one: say the Portuguese for the political party, then the Portuguese for the parliament.
\end{quote}

\subsection*{You'll want to know: eleger}
\textbf{eleger} — \textquotedblleft{}to elect\textquotedblright{}.

\begin{quote}
Elegeram uma mulher para presidente da câmara.
\end{quote}

\begin{quote}
\emph{They elected a woman as mayor.}
\end{quote}

Two participles: \emph{elegido} with \emph{ter}, \emph{eleito} with \emph{ser}: \emph{foi eleita}.

\begin{cousinweb}
\emph{eleger} is a learned borrowing from Latin \emph{ēligere} \textquotedblleft{}to pick out, to choose\textquotedblright{}, built from \emph{ex} and \emph{legere} \textquotedblleft{}to gather, to choose\textquotedblright{} (Lewis \& Short).
\end{cousinweb}

\begin{grammarlens}[title={What you've built}]
One more word for telling, explaining and describing at B1.
\end{grammarlens}

\subsection*{Your turn}
\begin{itemize}
\raggedright
  \item \emph{Say it:} \emph{eleger}
  \item \emph{Say it:} \emph{eleger}, once more
  \item \emph{Say it:} a sentence of your own with \emph{eleger}
  \item \emph{Recall:} say the Portuguese for the political party, then the Portuguese for the parliament, then say \emph{eleger} again
\end{itemize}

\begin{tcolorbox}[breakable,colback=teal!4,colframe=teal!35!black,title={Before you move on}]
What does eleger mean? (\textquotedblleft{}To elect\textquotedblright{}.) Say it once more.
\end{tcolorbox}
\section[o eleitor]{o eleitor — the voter}

\label{lesson:PT-C380-eleitor}



\begin{quote}
Before the new one: say the Portuguese for the parliament, then the Portuguese for to elect.
\end{quote}

\subsection*{You'll want to know: o eleitor}
\textbf{o eleitor} — \textquotedblleft{}the voter\textquotedblright{}.

\begin{quote}
Muitos eleitores não votaram este ano.
\end{quote}

\begin{quote}
\emph{Many voters did not vote this year.}
\end{quote}

Feminine \emph{a eleitora}.

\begin{cousinweb}
\emph{eleitor} comes from Latin \emph{ēlēctor} \textquotedblleft{}a chooser\textquotedblright{} (Lewis \& Short), on \emph{ēligere} \textquotedblleft{}to pick out\textquotedblright{}; the \emph{-it-} shows the popular change of \emph{-ct-}.
\end{cousinweb}

\begin{grammarlens}[title={What you've built}]
One more word for telling, explaining and describing at B1.
\end{grammarlens}

\subsection*{Your turn}
\begin{itemize}
\raggedright
  \item \emph{Say it:} \emph{o eleitor}
  \item \emph{Say it:} \emph{o eleitor}, once more
  \item \emph{Say it:} a sentence of your own with \emph{o eleitor}
  \item \emph{Recall:} say the Portuguese for the parliament, then the Portuguese for to elect, then say \emph{o eleitor} again
\end{itemize}

\begin{tcolorbox}[breakable,colback=teal!4,colframe=teal!35!black,title={Before you move on}]
What does o eleitor mean? (\textquotedblleft{}The voter\textquotedblright{}.) Say it once more.
\end{tcolorbox}
